% Vietnamese translation for OI Public Library
% Source-File: IOI中国国家候选队论文/2008/Day2/6.余林韵《运用化归思想解决信息学中的数列问题》/讲稿.doc
\documentclass[11pt]{article}
\usepackage{oipl-vn}

\newcommand{\Oh}{\mathcal{O}}

\title{Dùng tư tưởng quy về để giải các bài toán dãy số}
\author{Yu Linyun\\Trường Trung học số 1 Phúc Châu}
\date{}

\begin{document}
\maketitle

\origtitle{Solving Sequence Problems by Reduction}
\sourcepath{IOI 2008 / Day 2 / Yu Linyun / talk-script.doc}

\section{Mở đầu}

Xin chào mọi người, tôi là Yu Linyun từ Trường Trung học số 1 Phúc Châu. Chủ
đề tôi trình bày hôm nay là \emph{Dùng tư tưởng quy về để giải các bài toán
dãy số}.

Trong đời sống hằng ngày, chúng ta thường gặp rất nhiều số được sắp theo một
quy luật nhất định, chẳng hạn khẩu hiệu khi huấn luyện quân sự
\[
  1,1,1,2,1,1,1,1,2,1,\ldots
\]
hoặc các năm
\[
  1980,1981,1982,1983,1984,1985,1986,\ldots
\]
Những ví dụ như vậy đều là các dãy số được sắp theo một quy tắc nào đó.

Dãy số là một nhánh quan trọng của toán học. Từ nhỏ ta đã học về dãy số, từ
những bài tìm quy luật ở tiểu học đến chương riêng trong toán phổ thông; các
bài liên quan đến dãy số cũng thường xuất hiện trong các kỳ thi toán trung
học. Trong công việc và đời sống, dãy số vẫn có vai trò quan trọng, ví dụ tỉ
lệ vàng trong phân tích cổ phiếu có liên hệ đến dãy Fibonacci.

Trong vài năm gần đây, các bài toán dãy số cũng nhiều lần xuất hiện trong
các cuộc thi tin học. Bài trình bày chọn bốn bài khá đại diện, phân tích theo
đặc điểm riêng của từng bài và thử rút ra một số kỹ thuật.

\section{Tư tưởng quy về}

Để giải bài toán dãy số, một mặt người giải cần kiến thức toán học phong phú;
mặt khác cần tư duy khai phá, biết vận dụng tư tưởng quy về, mạnh dạn phỏng
đoán và liên tục đào sâu bản chất của bài toán.

``Quy về'' gồm chuyển hóa và quy kết. Khi giải một bài toán \(A\), ta thường
chuyển nó qua một quá trình nào đó thành bài toán \(B\) đã giải được hoặc dễ
giải hơn, rồi từ lời giải của \(B\) quay lại lời giải của \(A\). Đó là tư
tưởng cơ bản của phương pháp quy về. Bản chất của nó là biến vấn đề mới
thành tri thức cũ đã nắm được, rồi từ đó hiểu và giải quyết vấn đề mới.

Ba yếu tố của phương pháp quy về là:
\begin{itemize}
  \item \term{đối tượng quy về}: ta quy về cái gì;
  \item \term{mục tiêu quy về}: quy về đâu;
  \item \term{con đường quy về}: quy về bằng cách nào.
\end{itemize}

Khi dùng tư tưởng này cho bài toán dãy số, ta cần nắm chắc đặc tính của dãy,
chẳng hạn tính chu kỳ hoặc tính theo giai đoạn; những đặc tính ấy đều có thể
trở thành con đường quy về.

\section{Ví dụ: Dispute}

Xét bài ví dụ \emph{Dispute}. Bài toán cho một công thức truy hồi của một dãy
số và yêu cầu tính một hạng cụ thể. Quy mô dữ liệu là
\[
  0\le N\le 10^8.
\]
Vì dữ liệu quá lớn, rõ ràng không thể truy hồi đơn giản từ \(1\) đến \(N\).

Nút thắt của truy hồi trực tiếp nằm ở quy mô dữ liệu quá lớn. Vậy tại sao
không thử giảm quy mô xuống một phạm vi chấp nhận được rồi mới truy hồi? Từ
đó ta có thuật toán đầu tiên: chia toàn bộ dãy thành nhiều đoạn để xử lý.

Thông qua tiền xử lý, tính các hạng của dãy \(F\) và lưu lại mọi giá trị
\(F_{100000k}\). Khi cho \(N\), ta tìm hạng \(F_{100000k}\) gần \(N\) nhất,
rồi dùng truy hồi đơn giản để tính tiếp. Số bước truy hồi nhiều nhất là
\(100000\), nên độ phức tạp là \(\Oh(100000)\), rất hiệu quả trong phạm vi
này.

Nhưng nếu tiếp tục tăng quy mô lên \(10^{10}\), \(10^{11}\), thậm chí
\(10^{15}\), thuật toán đầu tiên không còn đảm đương được. Ta không thể tiền
xử lý nhiều như vậy. Giả sử \(10^7\) bước truy hồi mất \(1\) giây, thì với
dữ liệu cỡ \(10^{15}\), tiền xử lý cần \(10^8\) giây, tức hơn \(3\) năm.

\section{Khai thác cấu trúc của dãy}

Quay lại ý tưởng của thuật toán đầu tiên. Mục tiêu quy về là đưa quy mô về
mức có thể chịu được; con đường quy về là tiền xử lý rồi chia dữ liệu thành
các đoạn hữu hạn. Khi quy mô tăng, tiền xử lý trở thành rào cản không vượt
qua được. Nguyên nhân thất bại là: dãy số rất đẹp, nhưng thuật toán đầu tiên
chỉ thô bạo chia nó thành các phần. Liệu có thể lợi dụng chính đặc tính của
dãy để chia đoạn một cách ``dịu dàng'' hơn không?

Quan sát lại dãy: hàm \(g(x,y)\) có rất nhiều hạng và trông đáng sợ, nhưng
có hai đặc điểm.
\begin{enumerate}
  \item Trong mọi hạng, số mũ của \(y\) chỉ là \(0\) hoặc \(1\).
  \item Cuối công thức lấy modulo một số nguyên tố \(9973\).
\end{enumerate}

Từ đặc điểm thứ nhất, nếu xem \(x\) là hằng số thì sau khi biến đổi:
\[
  g(x,y)=((x^5-x+7)y+(-x^5+x^3+3x)) \bmod 9973.
\]
Vì thế \(g(x,y)\) là tuyến tính theo \(y\), và dãy \(F_n\) là một truy hồi
tuyến tính bậc một theo từng bước thích hợp. Liên hệ với việc tính hạng thứ
\(N\) của truy hồi tuyến tính hệ số hằng bằng nhân ma trận, ta điều chỉnh mục
tiêu quy về: tìm cách chia đoạn hợp lý để biến dãy thành truy hồi tuyến tính
hệ số hằng.

Đặc điểm thứ nhất cho hướng giải, còn đặc điểm thứ hai làm cho hướng ấy khả
thi. Đặt \(M=9973\). Phân tích thêm cho thấy có thể viết
\[
  F_{kM}=A F_{(k-1)M}+B.
\]
Do đó ta dùng nhân ma trận để nhanh chóng tính hạng \(F_{kM}\) gần \(N\)
nhất, rồi truy hồi thêm \(\Oh(M)\) bước để ra đáp án. Độ phức tạp toàn bộ là
\[
  \Oh\!\left(M+\log_2\frac{N}{M}\right).
\]

Vì thời gian trình bày có hạn, ba ví dụ còn lại không được chia sẻ chi tiết
ở đây; người nghe có thể tham khảo luận văn để thảo luận thêm.

\section{Kết luận}

Bài toán dãy số biến hóa muôn hình. Quy về không phải vạn năng, nhưng nếu ta
nắm chắc bản chất bài toán, quan sát cẩn thận, dám phỏng đoán, có tinh thần
sáng tạo và phát hiện được vẻ đẹp của dãy số, ta có thể vượt qua khó khăn để
đến lời giải. Điều này cũng đúng với các bài toán tin học khác.

Xin cảm ơn mọi người đã lắng nghe. Rất mong được trao đổi và trả lời câu hỏi.

\end{document}
